%%
%% Chapter 2: Related Work
%% Independent file, included by ../main.tex
%%

\section{Related Work}
\label{sec:rel}

\subsection{Sparse-View CT and CBCT Reconstruction}
\label{sec:rel:cbct}

Analytical reconstruction (FDK, filtered back-projection, ART-family
iterative solvers) degrades sharply when the number of projections is
small, producing severe streak artifacts. Classical regularization with
total variation and dictionary priors mitigates the artifacts but
struggles on anatomy-correlated noise and requires costly tuning.
Deep learning approaches first dominated the field as post-processing
networks that clean artifacts on FDK images, then progressed to
end-to-end and unrolled models that alternate data-consistency steps
with learned proximal operators, and later to deep generative models in
which diffusion or score-based samplers synthesize volumes from sparse
views \cite{song2021scorebased}. More recently, implicit neural
representations (INRs) and explicit radiance-field renderers have been
adopted for CT, treating the volume as parameters of a differentiable
function that is fit by projection consistency \cite{mildenhall2020nerf}.
The DIF-Gaussian method \cite{lin2024difgaussian} replaces the density
field of such renderers with 3D Gaussian primitives, dramatically
accelerating both training and inference while maintaining high
fidelity, and is the backbone adopted in this work.

\subsection{3D Gaussian Splatting in Medical Imaging}
\label{sec:rel:3dgs}

3D Gaussian splatting (3DGS) \cite{kerbl2023gaussiansplats} represents a
scene as a collection of small anisotropic Gaussians and renders it with
highly optimized differentiable rasterization. Originally designed for
novel-view synthesis from dense photographs, 3DGS has been carried into
medical imaging because its primitives are a natural fit for tomographic
volumes: opacity and scattering can be rendered along X-ray rays through
accumulated transparency, and the representation supports billion-voxel
domains without the memory cost of dense grids. Sparse-view CBCT
reconstruction with 3DGS achieves competitive quality at a fraction of
the training cost of NeRF-style field renderers \cite{lin2024difgaussian}.
Our work builds on this line by introducing a latent
stochastic layer on top of the splatting decoder.

\subsection{Variational Latent Models and Posterior Collapse}
\label{sec:rel:vae}

Variational autoencoders (VAEs) \cite{kingma2014vae} learn amortized
posteriors over latent variables by maximizing the evidence lower bound
(ELBO). Conditional VAEs (CVAEs) \cite{sohn2015cvae} condition both the
prior and the posterior on auxiliary inputs and are the standard tool
for structured prediction with latent stochasticity. A well-known
degeneracy of VAE training is \emph{posterior collapse}: the posterior
approximation collapses onto the prior, the KL term vanishes, and the
latent code becomes uninformative \cite{bowman2016generating,
lucas2019donot}. Common remedies include KL
annealing, free bits \cite{kingma2016freebits}, and architectural
changes that limit information leakage from the decoder to the latent
\cite{zhao2019infovae, dieng2019avoiding}. We observe the same
degeneracy in the splatting-based reconstruction setting and show that
the dominant cause is \emph{information leakage through the decoder
conditioning}: when the decoder views dense-view features during
training, the latent is redundant and collapses, while a latent that is
the \emph{only} channel for dense information retains its role.

\subsection{Neural Processes}
\label{sec:rel:np}

Neural processes (NPs) \cite{garnelo2018neural} combine the flexibility
of stochastic processes with the scalability of neural networks: a
context encoder produces a global representation of observations, which
conditions the prediction distribution through a latent variable. Neural
process families such as latent NPs \cite{garnelo2018latent} and
convolutional NPs \cite{gordon2019convolutional} provide tractable
variational objectives and permutation-invariant context encoding, and
have been extended to point clouds and medical data. Our dense-NP
formulation adapts this framework to tomographic reconstruction: sparse
views $C$ play the role of context, dense views $D$ (and, in the P5
variant, the ground-truth volume $Y$) condition the posterior during
training, and the variational prior computed from $C$ is the sole source
of latent stochasticity at deployment.

\subsection{Uncertainty Estimation in Image Reconstruction}
\label{sec:rel:unc}

Uncertainty in inverse problems is commonly categorized into aleatoric
(irreducible, sensory/observation noise) and epistemic (model and data
ambiguity). Sampling-based Bayesian deep learning estimates epistemic
uncertainty via ensembles \cite{lakshminarayanan2017deep} or
Monte-Carlo dropout \cite{gal2016dropout}; reconstruction-specific
methods include probabilistic unrolling and conditional normalizing
flows. Our approach obtains both flavors in a
single forward pass: latent Monte-Carlo sampling captures epistemic
ambiguity over scene content, while a homoscedastic observation model
quantifies measurement noise. Both are essential for clinical
decision-making, which is the motivation for the probabilistic
redesign in this paper.